% =========================================================================
% Chapter 3: Review of Literature
% =========================================================================

\chapter{Review of Literature}
\label{ch:lit_review}

The literature is arranged in three tiers: international work on antimicrobial resistance, plant-extract screening, and the genus \textit{Dipsacus}; national Indian work on plant extracts against multidrug-resistant clinical strains, with Himalayan taxa where those data exist; and regional work from Jammu and Kashmir. Entries are given as Author(s) (Year) and are ordered chronologically within each tier. Core findings, target organisms, extract media, and phytochemical classes are retained from the cited papers.

\section{International work}
\label{sec:lit_international}

\begin{itemize}

\item \citet{bauer1966}: described the standardised single-disc diffusion method for measuring antibiotic susceptibility as a circular zone of inhibition around a paper disc on an inoculated agar plate. The method remains the template for later plant-extract disc tests, although crude extracts have no CLSI interpretive breakpoints.

\item \citet{cowan1999}: reviewed plant products as antimicrobial agents and catalogued phenolics, terpenoids, alkaloids, flavonoids, tannins, and saponins as the main active classes. The review emphasised that plant extracts are chemically mixed and that activity depends on solvent polarity as well as on the test organism.

\item \citet{nikaido2003}: restated the molecular basis of Gram-negative outer-membrane permeability. Lipid~A, core oligosaccharide, and divalent-cation bridges exclude many hydrophobic molecules, while porins admit selected hydrophilic solutes. The paper is the standard account of why many plant phenolics fail against Gram-negative rods.

\item \citet{burt2004}: reviewed antibacterial essential oils and their constituents (terpenes and phenylpropanoids) against food-borne and clinical bacteria, including \textit{Escherichia coli} and \textit{Pseudomonas aeruginosa}. Lipophilic terpenoids were reported to swell membranes, raise proton leak, and collapse the proton-motive force.

\item \citet{cushnie2005}: reviewed antibacterial flavonoids. Activity was linked to B-ring hydroxylation and lipophilicity. Reported targets included DNA gyrase, membranes, and, in some studies, energy metabolism. Test organisms commonly included \textit{Staphylococcus aureus} and \textit{E.\ coli}.

\item \citet{valgas2007}: compared screening methods for natural products and showed that agar-well and disc-diffusion diameters are screening endpoints, not interchangeable with broth MIC values.

\item \citet{liebold2011}: tested lipophilic root extracts of \textit{Dipsacus sylvestris} Huds.\ against \textit{Borrelia burgdorferi} sensu stricto in vitro and reported growth inhibition. The extract medium was organic, not aqueous, and the target was a spirochaete rather than a Gram-negative rod.

\item \citet{zhao2011}: reviewed the phytochemistry of \textit{Dipsacus}. The dominant classes were triterpenoid saponins (oleanolic acid and hederagenin glycosides, including asperosaponin~VI / dipsacus saponin~C), iridoid glycosides (loganin, sweroside, cantleyoside, sylvestrosides), and phenolics. Most pharmacological records concerned \textit{D.\ asper}, not \textit{D.\ inermis}.

\item \citet{xie2015}: surveyed flavonoid structure--activity relationships against bacteria, including multidrug-resistant isolates, and again placed activity in membrane and enzyme effects rather than in a single clinical breakpoint.

\item \citet{blair2015}: reviewed molecular mechanisms of antibiotic resistance, including $\beta$-lactamases, porin loss, and RND efflux (AcrAB-TolC; MexAB-OprM). The review supplies the international mechanistic context for why Gram-negative clinical isolates are hard to inhibit with crude plant extracts.

\item \citet{li2015}: reviewed efflux-mediated resistance in Gram-negative bacteria, with detailed accounts of RND pumps in \textit{E.\ coli}, \textit{K.\ pneumoniae}, and \textit{P.\ aeruginosa}.

\item \citet{apgiv2016}: placed Dipsacaceae within an expanded Caprifoliaceae. Teasels, including \textit{Dipsacus}, are therefore cited here as Caprifoliaceae sensu lato.

\item \citet{balouiri2016}: reviewed agar disc diffusion, well diffusion, and dilution MIC methods for plant extracts. Zone diameter was defined as a screening measure; MIC remains the reference quantitative endpoint.

\item \citet{vanvuuren2019}: listed pitfalls of plant-extract zone-of-inhibition studies, including non-standard media, unstated inocula, and the absence of clinical breakpoints for crude extracts.

\item \citet{yu2019}: isolated triterpenoids and triterpenoid saponins from \textit{Dipsacus asper} roots and tested cytotoxic and antibacterial activities of the purified compounds. The work confirms antibacterial interest in teasel saponins but used isolated molecules, not a crude aqueous extract of \textit{D.\ inermis}.

\item \citet{oszmianski2020}: analysed iridoids and phenolics in leaves and roots of \textit{Dipsacus fullonum} L.\ by UPLC-PDA-MS/MS (loganic acid, loganin, sweroside, cantleyoside, sylvestroside~III; phenolic acids and flavones). Root extract produced inhibition zones against \textit{Staphylococcus aureus} DSM~799 and \textit{E.\ coli} ATCC~10536 in a well-diffusion test; other bacteria in that panel, including \textit{P.\ aeruginosa}, were not inhibited under the same conditions.

\item \citet{jubair2021}: reviewed plant-derived compounds against multidrug-resistant bacteria and contrasted multi-target plant chemistry with single-target synthetic antibiotics.

\item \citet{amrcollaborators2022}: in the GRAM analysis, estimated 4.95 million deaths associated with bacterial AMR in 2019 and 1.27 million deaths attributable to resistant infection, with a high burden in South Asia.

\item \citet{skala2023}: reviewed \textit{Dipsacus} iridoids, saponins, and phenolics and their reported anti-inflammatory and antimicrobial effects. Water-soluble iridoid glycosides were noted as recoverable in aqueous decoctions.

\item \citet{who2024}: issued the Bacterial Priority Pathogens List. Carbapenem-resistant Enterobacterales (including \textit{K.\ pneumoniae} and \textit{E.\ coli}) and carbapenem-resistant \textit{P.\ aeruginosa} remain in the critical group.

\item \citet{drakhshaan2026}: studied whole-plant \textit{Dipsacus inermis} from Muzaffarabad (northern Pakistan) as dichloromethane and methanol extracts. GC-MS of the dichloromethane fraction was dominated by terpenoids, steroids, and sesquiterpenoids; HPLC of the methanol fraction detected polar phenolics including rutin and ferulic acid. At 100\,\textmu g\,mL$^{-1}$ in well diffusion, dichloromethane zones were $15.93 \pm 0.11$\,mm against \textit{E.\ coli} and $16.83 \pm 0.29$\,mm against \textit{P.\ aeruginosa}; methanol values were $15.80 \pm 0.34$\,mm and $15.17 \pm 0.29$\,mm. Broth MIC values ranged from 1.562 to 6.25\,mg\,mL$^{-1}$. \textit{Klebsiella pneumoniae} was not tested, and no purely aqueous extract was reported.

\end{itemize}

\section{National work (India)}
\label{sec:lit_national}

\begin{itemize}

\item \citet{ahmadbeg2001}: screened alcoholic extracts of 45 Indian medicinal plants against multidrug-resistant clinical isolates of \textit{S.\ aureus}, \textit{E.\ coli}, \textit{Salmonella paratyphi}, \textit{Shigella dysenteriae}, \textit{Bacillus subtilis}, and \textit{Candida albicans}. Twelve plants, including \textit{Lawsonia inermis}, \textit{Holarrhena antidysenterica}, \textit{Terminalia chebula}, \textit{Terminalia bellirica}, \textit{Camellia sinensis}, and \textit{Punica granatum}, showed marked activity. Phytochemical tests on active extracts detected alkaloids, glycosides, flavonoids, phenolics, and saponins.

\item \citet{aqil2005}: tested Indian plant extracts against clinical $\beta$-lactamase-producing MRSA and MSSA. Extracts of \textit{Camellia sinensis}, \textit{Delonix regia}, \textit{Holarrhena antidysenterica}, \textit{Lawsonia inermis}, \textit{Punica granatum}, \textit{Terminalia chebula}, and \textit{Terminalia bellirica} produced zones of 11--27\,mm. MIC values by tube dilution ranged from 1.3 to 8.2\,mg\,mL$^{-1}$. TLC indicated alkaloids, glycosides, flavonoids, phenols, and saponins. Synergy with tetracycline (and, for tea, ampicillin) was reported for several extracts.

\item \citet{ahmadaqil2007}: tested bioactive extracts of 15 Indian medicinal plants against ES$\beta$L-producing multidrug-resistant enteric bacteria. Ethanolic extracts inhibited ES$\beta$L isolates; MIC values against that panel ranged from 0.31 to 6.25\,mg\,mL$^{-1}$. \textit{Plumbago zeylanica}, \textit{Hemidesmus indicus}, \textit{Acorus calamus}, and \textit{Camellia sinensis} ranked among the more active plants. The paper is a national reference for plant extracts against Gram-negative clinical strains, not a Himalayan floristic survey.

\item \citet{kumar2013}: reviewed flavonoids as plant secondary metabolites with reported antibacterial, antioxidant, and anti-inflammatory actions, providing the Indian review context for flavonoid-positive extracts.

\item \citet{tao2020}: reviewed \textit{Dipsacus asper} (widely used in Asian traditional medicine, including Indian herbal trade as a related teasel) for saponins and iridoids. Direct Indian antibacterial trials of \textit{D.\ inermis} remain scarce in indexed sources.

\item \citet{adhikari2023}: worked on needles of \textit{Taxus wallichiana} Zucc.\ (Himalayan yew) at the G.B.\ Pant National Institute of Himalayan Environment. Bioassay-guided fractionation, TLC-bioautography, GC-MS/LC-MS, FTIR, and NMR identified fatty acids, alkaloids, and related metabolites with activity against bacteria, actinobacteria, and fungi. The paper is an indexed Himalayan Indian antimicrobial study on a high-altitude gymnosperm, not on \textit{Dipsacus}.

\item \citet{bhardwaj2023}: tested methanol extracts and fractions of leaves and fruit of \textit{Taxus wallichiana} by agar-well diffusion against five bacteria, with ciprofloxacin as standard. Leaf and fruit methanol extracts produced maximum zones of $18.0 \pm 0.0$\,mm against \textit{Bacillus subtilis} and $19 \pm 0.2$\,mm against \textit{S.\ aureus}. Leaf extract was positive for glycosides, flavonoids, phenol, tannins, and terpenoids; fruit extract for flavonoids, saponins, and terpenoids.

\item \citet{sati2025}: described methods behind the WHO 2024 Bacterial Priority Pathogens List. Carbapenem-resistant Enterobacterales and \textit{P.\ aeruginosa} remain global and Indian research targets.

\end{itemize}

\section{Regional work (Jammu and Kashmir)}
\label{sec:lit_regional}

\begin{itemize}

\item \citet{guna2018}: collected eleven ethnomedicinal plants from the Kashmir Valley and prepared hexane and methanol extracts by cold percolation. Paper-disc diffusion was run against \textit{E.\ coli}, \textit{S.\ aureus}, \textit{B.\ subtilis}, \textit{P.\ aeruginosa}, \textit{Proteus vulgaris}, and \textit{Candida albicans} at 1.25, 2.5, and 5\,mg/disc. All plants inhibited at least one organism. \textit{Nepeta cataria}, \textit{Allium consangium}, \textit{Artemisia} (\textit{biensis}), and \textit{Ocimum sanctum} were the most active. Methanol extracts were generally stronger than hexane. Bacteria were more often inhibited than \textit{C.\ albicans}. \textit{Dipsacus} was not in the panel.

\item \citet{hassan2020}: tested \textit{Dipsacus inermis} Wall.\ collected in Kashmir for anti-inflammatory activity. In LPS-stimulated macrophages the extract reduced nitric oxide, iNOS, and COX-2 and limited NF-$\kappa$B signalling (I$\kappa$B$\alpha$ phosphorylation and p65 nuclear translocation). In mice with carrageenan paw oedema the extract reduced swelling. The paper documents a Kashmir collection of the same species used here, but it is not an antibacterial disc-diffusion study.

\item \citet{nabi2022}: prepared $n$-hexane, ethyl acetate, and methanol root extracts of \textit{Skimmia anquetilia} N.P.\ Taylor and Airy Shaw from the Kashmir Himalaya. Disc or related diffusion tests were run against MDR strains \textit{P.\ aeruginosa} MTCC~424, \textit{E.\ coli} MTCC~739, \textit{K.\ pneumoniae} MTCC~139, \textit{Salmonella} Typhi MTCC~3224, and \textit{S.\ aureus} MTCC~96. Ethyl acetate gave the largest zones: 18\,mm against \textit{P.\ aeruginosa} and 17\,mm against \textit{S.\ aureus}. MIC of the ethyl acetate extract against \textit{S.\ aureus} was 4\,mg\,mL$^{-1}$. FTIR and GC-MS showed mixed aliphatics, aromatics, and fatty-acid derivatives. The study is a University of Kashmir report on a temperate Himalayan shrub against the same three Gram-negative genera used in this dissertation.

\item \citet{ji2026}: prepared ZnO nanoparticles from an aqueous extract of \textit{D.\ inermis} collected in Kupwara, Kashmir, and reported activity against \textit{S.\ aureus} and \textit{E.\ coli}. The assay was on nanoparticles, not on a crude aqueous extract against \textit{K.\ pneumoniae}.

\item \citet{drakhshaan2026}: although the plant was collected in northern Pakistan rather than in the Union Territory of Jammu and Kashmir, it is the closest indexed antibacterial data set on \textit{D.\ inermis} itself and is therefore noted again here as the regional Himalayan comparator. It used dichloromethane and methanol, not water, and omitted \textit{K.\ pneumoniae}.

\end{itemize}

\section{Synthesis}
\label{sec:lit_synthesis}

Indexed international work establishes AMR as a current clinical problem, disc diffusion as a screening method, and \textit{Dipsacus} as a saponin- and iridoid-rich genus with scattered antibacterial reports, mostly on organic extracts or on species other than \textit{D.\ inermis}. Indian work shows that crude plant extracts can inhibit MDR staphylococci and ES$\beta$L enteric bacteria, and Himalayan taxa such as \textit{Taxus wallichiana} have documented zones against Gram-positive cocci. Kashmir studies have tested temperate herbs and \textit{Skimmia anquetilia} against \textit{E.\ coli}, \textit{P.\ aeruginosa}, and \textit{K.\ pneumoniae}, and have examined \textit{D.\ inermis} for inflammation and as a nanoparticle reducing agent. What remains thinly documented is a simple aqueous extract of Kashmir \textit{D.\ inermis}, screened qualitatively for phytochemicals and tested by disc diffusion against those three clinical pathogenic stocks. That is the gap addressed in the following chapters.
